% Toán tư duy lớp 4 – Chủ đề 1 (bổ sung): mỗi dạng đủ 4 câu (1 ví dụ mẫu + 3 bài tự luyện) cho Luyện mỗi ngày.
% Em soạn, đã tự giải lại và kiểm tra phép tính bằng Python (2026-11-18).

\begin{bt}%[4D1H1-1]
% nguon: Toán tư duy lớp 4 – Bài 1 SGK Toán 4 KNTT
	Nếu thay chữ số $5$ của số $45\,678$ bằng chữ số $9$ thì được số mới. Số mới hơn số cũ bao nhiêu đơn vị?
	\loigiai{
		Chữ số $5$ ở hàng nghìn, có giá trị $5\,000$. Thay bằng $9$ thì có giá trị $9\,000$.\\
		Số mới là $49\,678$, hơn số cũ: $9\,000 - 5\,000 = 4\,000$ (đơn vị).\\
		Thử lại: $49\,678 - 45\,678 = 4\,000$.
	}
\end{bt}

\begin{bt}%[4D1H2-1]
% nguon: Toán tư duy lớp 4 – Bài 2 SGK Toán 4 KNTT
	Tính bằng cách thuận tiện: $4 \times 125 \times 25 \times 8$.
	\loigiai{
		Ghép các cặp nhân ra số tròn: $4 \times 25 = 100$ và $125 \times 8 = 1\,000$.\\
		$4 \times 125 \times 25 \times 8 = (4 \times 25) \times (125 \times 8) = 100 \times 1\,000 = 100\,000$.
	}
\end{bt}

\begin{bt}%[4D1V2-2]
% nguon: Toán tư duy lớp 4 – Bài 2 SGK Toán 4 KNTT
	Tìm số có hai chữ số $\overline{ab}$, biết: $\overline{ab} \times 5 = \overline{1b5}$.
	\loigiai{
		$\overline{1b5}$ nằm trong khoảng từ $100$ đến $199$, nên $\overline{ab}$ từ $20$ đến $39$ ($20 \times 5 = 100$, $40 \times 5 = 200$).\\
		$\overline{1b5}$ có tận cùng là $5$ nên $\overline{ab}$ là số lẻ: $21, 23, \dots, 39$.\\
		Ta cần chữ số hàng chục của tích đúng bằng chữ số $b$ (hàng đơn vị của $\overline{ab}$):\\
		$21 \times 5 = 105$, $23 \times 5 = 115$, $25 \times 5 = 125$, $27 \times 5 = 135$, $29 \times 5 = 145$: không số nào được.\\
		$31 \times 5 = 155$, $33 \times 5 = 165$, $35 \times 5 = 175$, $37 \times 5 = 185$: không được; $39 \times 5 = 195$: được!\\
		Vậy $\overline{ab} = 39$.
	}
\end{bt}

\begin{bt}%[4D1H2-3]
% nguon: Toán tư duy lớp 4 – Bài 2 SGK Toán 4 KNTT
	Điền dấu $\times$ hoặc $:$ vào mỗi ô trống để được phép tính đúng: $48 \ \square\ 6 \ \square\ 2 = 16$.
	\loigiai{
		Thử các cách điền (tính từ trái sang phải):\\
		$48 : 6 \times 2 = 8 \times 2 = 16$. Đúng!\\
		$48 \times 6 : 2 = 144$; \quad $48 : 6 : 2 = 4$; \quad $48 \times 6 \times 2 = 576$.\\
		Vậy phép tính đúng là $48 : 6 \times 2 = 16$.
	}
\end{bt}

\begin{bt}%[4D1N3-1]
% nguon: Toán tư duy lớp 4 – Bài 3 SGK Toán 4 KNTT
	Tìm số chẵn lớn nhất có bốn chữ số và số lẻ nhỏ nhất có bốn chữ số.
	\loigiai{
		Số lớn nhất có bốn chữ số là $9\,999$ (số lẻ), nên số chẵn lớn nhất có bốn chữ số là $9\,998$.\\
		Số nhỏ nhất có bốn chữ số là $1\,000$ (số chẵn), nên số lẻ nhỏ nhất có bốn chữ số là $1\,001$.
	}
\end{bt}

\begin{bt}%[4D1H3-2]
% nguon: Toán tư duy lớp 4 – Bài 3 SGK Toán 4 KNTT
	Hiệu của hai số tự nhiên là $37$. Hai số đó có thể cùng là số chẵn không? Có thể cùng là số lẻ không? Vì sao?
	\loigiai{
		Chẵn $-$ chẵn $=$ chẵn; lẻ $-$ lẻ $=$ chẵn.\\
		Hiệu là $37$, một số lẻ, nên hai số không thể cùng chẵn, cũng không thể cùng lẻ.\\
		Vậy hai số đó phải có một số chẵn và một số lẻ.
	}
\end{bt}

\begin{bt}%[4D1N3-3]
% nguon: Toán tư duy lớp 4 – Bài 3 SGK Toán 4 KNTT
	Lớp 4A có $25$ học sinh. Cô giáo xếp mỗi bàn $2$ bạn. Có bàn nào chỉ có $1$ bạn ngồi không? Cần ít nhất bao nhiêu bàn?
	\loigiai{
		$25$ là số lẻ nên không chia hết thành từng cặp $2$ bạn: $25 : 2 = 12$ (dư $1$).\\
		Xếp được $12$ bàn mỗi bàn $2$ bạn, còn $1$ bạn ngồi một bàn riêng.\\
		Vậy có $1$ bàn chỉ có $1$ bạn và cần ít nhất $12 + 1 = 13$ (bàn).
	}
\end{bt}

\begin{bt}%[4D1V3-3]
% nguon: Toán tư duy lớp 4 – Bài 3 SGK Toán 4 KNTT
	Có $9$ tấm thẻ ghi các số $1, 2, 3, 4, 5, 6, 7, 8, 9$. Có thể chia $9$ tấm thẻ thành hai nhóm sao cho tổng các số ở hai nhóm bằng nhau không? Vì sao?
	\loigiai{
		Tổng các số trên $9$ tấm thẻ: $1 + 2 + \dots + 9 = (1 + 9) \times 9 : 2 = 45$.\\
		Nếu hai nhóm có tổng bằng nhau thì tổng chung phải bằng hai lần tổng một nhóm, tức là số chẵn.\\
		$45$ là số lẻ, nên không thể chia được.
	}
\end{bt}

\begin{bt}%[4D1H4-1]
% nguon: Toán tư duy lớp 4 – Bài 4 SGK Toán 4 KNTT
	Tính giá trị của biểu thức $a \times b - c : 4$ với $a = 125$, $b = 8$, $c = 400$.
	\loigiai{
		Thay số: $125 \times 8 - 400 : 4$.\\
		Làm nhân, chia trước: $125 \times 8 = 1\,000$; $400 : 4 = 100$.\\
		Rồi trừ: $1\,000 - 100 = 900$.
	}
\end{bt}

\begin{bt}%[4D1H4-3]
% nguon: Toán tư duy lớp 4 – Bài 4 SGK Toán 4 KNTT
	Một hình vuông có chu vi là $P$. Viết biểu thức tính độ dài cạnh hình vuông theo $P$. Tính cạnh và diện tích hình vuông khi $P = 36$ cm.
	\loigiai{
		Chu vi bằng cạnh nhân $4$, nên cạnh $= P : 4$.\\
		Khi $P = 36$ cm: cạnh là $36 : 4 = 9$ (cm).\\
		Diện tích: $9 \times 9 = 81$ (cm$^2$).
	}
\end{bt}

\begin{bt}%[4D1H5-1]
% nguon: Toán tư duy lớp 4 – Bài 5 SGK Toán 4 KNTT
	Một xe tải chở gạo $3$ chuyến, mỗi chuyến $1\,250$ kg. Số gạo đó được chia đều cho $25$ cửa hàng. Hỏi mỗi cửa hàng nhận được bao nhiêu ki-lô-gam gạo?
	\loigiai{
		Bước 1. Số gạo xe chở được: $1\,250 \times 3 = 3\,750$ (kg).\\
		Bước 2. Mỗi cửa hàng nhận: $3\,750 : 25 = 150$ (kg).\\
		Đáp số: $150$ kg gạo.
	}
\end{bt}

\begin{bt}%[4D1H5-2]
% nguon: Toán tư duy lớp 4 – Bài 5 SGK Toán 4 KNTT
	Tìm một số, biết rằng lấy số đó nhân với $4$ rồi trừ đi $36$ thì được $100$.
	\loigiai{
		Tính ngược từ cuối (trừ thành cộng, nhân thành chia):\\
		Trước khi trừ $36$: $100 + 36 = 136$.\\
		Trước khi nhân $4$: $136 : 4 = 34$.\\
		Số cần tìm là $34$. Thử lại: $34 \times 4 - 36 = 136 - 36 = 100$.
	}
\end{bt}

\begin{bt}%[4D1V5-3]
% nguon: Toán tư duy lớp 4 – Bài 5 SGK Toán 4 KNTT
	Lớp em có $32$ bạn đi tham quan bằng $8$ chiếc thuyền gồm thuyền lớn và thuyền nhỏ. Mỗi thuyền lớn chở $5$ bạn, mỗi thuyền nhỏ chở $3$ bạn, các thuyền đều chở vừa đủ. Hỏi có mấy thuyền mỗi loại?
	\loigiai{
		Giả sử cả $8$ thuyền đều là thuyền nhỏ thì chở được $3 \times 8 = 24$ (bạn).\\
		Còn thiếu chỗ cho: $32 - 24 = 8$ (bạn).\\
		Mỗi thuyền lớn chở hơn thuyền nhỏ $5 - 3 = 2$ (bạn), nên số thuyền lớn: $8 : 2 = 4$ (thuyền).\\
		Số thuyền nhỏ: $8 - 4 = 4$ (thuyền).\\
		Thử lại: $4 \times 5 + 4 \times 3 = 20 + 12 = 32$ (bạn).
	}
\end{bt}
